% Problems Folder Detailed API Documentation (Current Implementation)
% To be included in master LaTeX document
%
% Usage: % Problems Folder Detailed API Documentation (Current Implementation)
% To be included in master LaTeX document
%
% Usage: % Problems Folder Detailed API Documentation (Current Implementation)
% To be included in master LaTeX document
%
% Usage: % Problems Folder Detailed API Documentation (Current Implementation)
% To be included in master LaTeX document
%
% Usage: \input{api/problems_folder_detailed_api}

\section{Problems Folder Detailed API Reference}
\label{sec:problems_folder_detailed_api}

This section provides comprehensive documentation for the actual problem modules implemented in the BioNetFlux problems folder, based on the current source code analysis.

\subsection{Problems Folder Structure}

The problems folder contains the following implemented modules:

\begin{itemize}
    \item \textbf{custom\_problem\_template.py}: Template for creating new problems
    \item \textbf{ooc\_problem.py}: OrganOnChip grid problem with configuration system
    \item \textbf{ooc\_config\_manager.py}: OrganOnChip-specific configuration manager
    \item \textbf{ks\_config\_manager.py}: Keller-Segel-specific configuration manager
\end{itemize}

\subsection{Configuration-Based Problem Architecture}
\label{subsec:config_architecture}

All current problem modules follow a configuration-driven architecture using TOML files and specialized configuration managers.

\subsubsection{Configuration Manager Pattern}

\begin{lstlisting}[language=Python, caption=Configuration Manager Pattern]
# Each problem type has its own configuration manager
from bionetflux.problems.ooc_config_manager import OoCConfigManager
from bionetflux.problems.ks_config_manager import KSConfigManager

def create_global_framework(config_file: Optional[str] = None):
    # Load configuration with validation
    config_manager = OoCConfigManager()  # or KSConfigManager()
    config = config_manager.load_config(config_file)
    
    # Extract resolved functions from config
    initial_conditions = config['initial_conditions']  # Already callables
    force_functions = config['force_functions']        # Already callables
    # ...existing code...
\end{lstlisting}

\subsection{OrganOnChip Configuration Manager}
\label{subsec:ooc_config_manager}

\subsubsection{OoCConfigManager Class}

\begin{lstlisting}[language=Python, caption=OoC Configuration Manager]
class OoCConfigManager(BaseConfigManager):
    def __init__(self):
        super().__init__("ooc")  # Config manager type
    
    def _setup_defaults(self):
        self.defaults = {
            'problem': {
                'name': 'OoC_Grid_Problem',
                'neq': 4,
                'problem_type': 'ooc'  # Config validation type
            },
            'physical_parameters': {
                'viscosity': {'nu': 1.0, 'mu': 2.0, 'epsilon': 1.0, 'sigma': 1.0},
                'reaction': {'a': 0.0, 'c': 0.0},
                'coupling': {'b': 1.0, 'd': 1.0, 'chi': 1.0}
            },
            'discretization': {
                'n_elements': 20,
                'tau': [0.5, 0.5, 0.5, 0.5]
            },
            'initial_conditions': {
                'u': 'zeros', 'omega': 'zeros', 'v': 'zeros', 'phi': 'zeros'
            },
            'force_functions': {
                'u': 'zeros', 'omega': 'zeros', 'v': 'zeros', 'phi': 'zeros'
            },
            'domain_initial_conditions': {},  # Domain-specific overrides
            'domain_force_functions': {}      # Domain-specific overrides
        }
\end{lstlisting}

\subsubsection{OoC Function Library}

The OoC configuration manager provides specialized functions:

\begin{lstlisting}[language=Python, caption=OoC Function Library]
def _setup_function_library(self):
    ooc_functions = {
        # OrganOnChip specific functions
        'ooc_viscous_profile': lambda s, t=0: s * (1 - s),
        'ooc_parabolic_flow': lambda s, t=0: 4 * s * (1 - s),
        'ooc_inlet_profile': lambda s, t=0: np.where(np.abs(s) < 0.1, 1.0, 0.0),
        
        # Time-dependent functions
        'sin_t': lambda s, t=0: np.sin(t) * np.ones_like(s),
        'cos_t': lambda s, t=0: np.cos(t) * np.ones_like(s),
        'exp_t': lambda s, t=0: np.exp(-t) * np.ones_like(s),
        
        # Combined space-time functions
        'traveling_wave': lambda s, t=0: np.sin(2 * np.pi * (s - t)),
        'diffusive_wave': lambda s, t=0: np.exp(-t) * np.sin(np.pi * s),
    }
\end{lstlisting}

\subsection{Keller-Segel Configuration Manager}
\label{subsec:ks_config_manager}

\subsubsection{KSConfigManager Class}

\begin{lstlisting}[language=Python, caption=KS Configuration Manager]
class KSConfigManager(BaseConfigManager):
    def __init__(self):
        super().__init__("ks")  # Config manager type
    
    def _setup_defaults(self):
        self.defaults = {
            'problem': {
                'name': 'KS_Traveling_Wave_Problem',
                'neq': 2,
                'problem_type': 'ks'
            },
            'physical_parameters': {
                'diffusion': {'mu': 2.0, 'nu': 1.0},
                'reaction': {'a': 0.0, 'b': 1.0},
                'chemotaxis': {
                    'chi': 'constant',    # Function name
                    'dchi': 'zeros'       # Function name
                }
            },
            'discretization': {
                'n_elements': 10,
                'tau': [1.0, 1.0]  # 2 equations
            },
            'initial_conditions': {
                'u': 'zeros',    # Cell density
                'phi': 'zeros'   # Chemical concentration
            }
        }
\end{lstlisting}

\subsubsection{KS Function Library}

\begin{lstlisting}[language=Python, caption=KS Function Library]
def _setup_function_library(self):
    ks_functions = {
        # Keller-Segel specific functions
        'ks_gaussian_blob': lambda s, t=0: np.exp(-(s - 2.0)**2),
        'ks_step_function': lambda s, t=0: np.where(np.abs(s - 2.0) < 0.5, 1.0, 0.0),
        
        # Chemotactic patterns
        'chemotactic_gradient': lambda s, t=0: s * np.exp(-s),
        'attraction_source': lambda s, t=0: np.exp(-((s - 2.0)**2) / 0.5),
    }
\end{lstlisting}

\subsection{Custom Problem Template}
\label{subsec:custom_problem_template}

The \texttt{custom\_problem\_template.py} provides a structured template for creating new problems.

\subsubsection{Template Structure}

\begin{lstlisting}[language=Python, caption=Custom Problem Template Structure]
def create_global_framework():
    """
    Custom problem template following MATLAB TestProblem.m structure.
    Returns: (problems, global_discretization, constraint_manager, problem_name)
    """
    
    # SECTION 1: PHYSICAL PARAMETERS (From MATLAB TestProblem.m)
    neq = 4  # 4-equation OrganOnChip system
    
    # Physical parameters following MATLAB TestProblem.m exactly
    nu, mu, epsilon, sigma = 1.0, 2.0, 1.0, 1.0  # Viscosity
    a, c = 0.0, 0.0                               # Reaction
    b, d, chi = 1.0, 1.0, 1.0                     # Coupling
    parameters = np.array([nu, mu, epsilon, sigma, a, b, c, d, chi])
    
    # SECTION 2: MATHEMATICAL FUNCTIONS
    def lambda_func(x): return np.ones_like(x)    # Constant function
    def dlambda_func(x): return np.zeros_like(x)  # Derivative
    
    # Initial conditions matching MATLAB TestProblem.m
    def initial_u(s, t=0.0): return np.sin(2 * np.pi * s)  # Non-trivial u
    def initial_omega(s, t=0.0): return np.zeros_like(s)    # Zero omega
    def initial_v(s, t=0.0): return np.zeros_like(s)        # Zero v  
    def initial_phi(s, t=0.0): return np.zeros_like(s)      # Zero phi
    
    # Source terms (all zero matching MATLAB)
    def source_u(s, t): return np.zeros_like(s)
    # ... similar for omega, v, phi
    
    # SECTION 3: GEOMETRY SETUP - ⚠️ CUSTOMIZE THIS SECTION ⚠️
    geometry = DomainGeometry("custom_geometry")
    # User adds geometry.add_domain() calls here
    
    # SECTION 4: PROBLEM CREATION (Automated from geometry)
    # SECTION 5: CONSTRAINT SETUP - ⚠️ CUSTOMIZE THIS SECTION ⚠️  
    # SECTION 6-7: GLOBAL DISCRETIZATION AND VALIDATION
\end{lstlisting}

\subsection{OrganOnChip Grid Problem}
\label{subsec:ooc_grid_problem}

The \texttt{ooc\_problem.py} implements a complete OrganOnChip problem with grid geometry.

\subsubsection{Main Framework Method}

\begin{lstlisting}[language=Python, caption=OoC Problem Framework]
def create_global_framework(geometry: Optional[DomainGeometry] = None,
                          config_file: Optional[str] = None):
    """
    OrganOnChip problem with grid geometry and configuration system.
    
    Args:
        geometry: Optional pre-defined geometry. Uses build_default_geometry() if None
        config_file: Optional TOML configuration file
    
    Returns:
        Tuple: (problems, global_discretization, constraint_manager, problem_name)
    """
    
    CONFIG_VALIDATION_TYPE = "ooc"      # For config file validation
    PROBLEM_TYPE = "organ_on_chip"      # For StaticCondensationFactory
    
    # Configuration loading with validation
    config_manager = OoCConfigManager()
    config = config_manager.load_config(config_file)
    
    # Extract configuration sections
    problem_config = config['problem']
    phys_params = config['physical_parameters']
    initial_conditions = config['initial_conditions']  # Already callables
    domain_initial_conditions = config.get('domain_initial_conditions', {})
    domain_force_functions = config.get('domain_force_functions', {})
    
    # Create problems from geometry
    problems = []
    for domain_id in range(geometry.num_domains()):
        domain_info = geometry.get_domain(domain_id)
        problem = Problem(neq, domain_info.domain_start, domain_info.domain_length,
                         parameters, "my_problem_type", name)
        # Set functions...
        problems.append(problem)
    
    # Setup constraints from geometry
    constraint_manager = setup_constraints_from_geometry(geometry, problems, neq)
    
    # Return framework components
    return problems, global_discretization, constraint_manager, problem_name
\end{lstlisting}

\subsubsection{Geometry Integration}

\begin{lstlisting}[language=Python, caption=Geometry Integration]
def build_default_geometry():
    """Build the default OoC grid geometry."""
    return build_grid_geometry()  # From geometry module

def setup_constraints_from_geometry(geometry: DomainGeometry, problems, neq: int):
    """Generate constraint manager from geometry connections."""
    constraint_manager = ConstraintManager()
    
    # Process boundary connections
    boundary_connections = geometry.get_boundary_connections()
    for conn in boundary_connections:
        if conn.is_exterior_boundary():
            for eq_idx in range(neq):
                constraint_manager.add_neumann(
                    equation_index=eq_idx,
                    domain_index=conn.domain1_id,
                    position=conn.parameter1,
                    data_function=lambda t: 0.0  # Homogeneous Neumann
                )
    
    # Process interior connections
    interior_connections = geometry.get_interior_connections()
    for conn in interior_connections:
        for eq_idx in range(neq):
            constraint_manager.add_trace_continuity(
                equation_index=eq_idx,
                domain1_index=conn.domain1_id,
                domain2_index=conn.domain2_id,
                position1=conn.parameter1,
                position2=conn.parameter2
            )
    
    return constraint_manager
\end{lstlisting}

\subsubsection{Domain-Specific Overrides}

\begin{lstlisting}[language=Python, caption=Domain-Specific Configuration]
# Apply domain-specific initial condition overrides
equation_names = ['u', 'omega', 'v', 'phi']

for domain_eq_key, func_name in domain_initial_conditions.items():
    # Parse "domain_<id>_<equation_name>"
    if domain_eq_key.startswith('domain_') and '_' in domain_eq_key[7:]:
        parts = domain_eq_key[7:].split('_', 1)
        domain_id = int(parts[0])
        equation_name = parts[1]
        
        if domain_id < len(problems) and equation_name in equation_names:
            equation_id = equation_names.index(equation_name)
            resolved_func = config_manager.function_resolver.resolve_function(func_name)
            problems[domain_id].set_initial_condition(equation_id, resolved_func)
            print(f"✓ Set initial condition: domain {domain_id}, equation '{equation_name}' -> {func_name}")
\end{lstlisting}

\subsection{Parameter Management}
\label{subsec:parameter_management}

\subsubsection{OrganOnChip Parameter Structure}

\begin{longtable}{|p{3cm}|p{2cm}|p{2cm}|p{6cm}|}
\hline
\textbf{Parameter} & \textbf{Index} & \textbf{Default} & \textbf{Physical Meaning} \\
\hline
\endhead

\texttt{nu} & 0 & 1.0 & Cell diffusion coefficient \\
\hline

\texttt{mu} & 1 & 2.0 & Potential diffusion coefficient \\
\hline

\texttt{epsilon} & 2 & 1.0 & Chemical diffusion coefficient \\
\hline

\texttt{sigma} & 3 & 1.0 & Velocity diffusion coefficient \\
\hline

\texttt{a} & 4 & 0.0 & Growth/decay rate \\
\hline

\texttt{b} & 5 & 1.0 & Velocity-potential coupling \\
\hline

\texttt{c} & 6 & 0.0 & Chemical reaction rate \\
\hline

\texttt{d} & 7 & 1.0 & Cell-chemical coupling \\
\hline

\texttt{chi} & 8 & 1.0 & Chemotactic sensitivity \\
\hline

\end{longtable}

\subsubsection{Keller-Segel Parameter Structure}

\begin{longtable}{|p{3cm}|p{2cm}|p{2cm}|p{6cm}|}
\hline
\textbf{Parameter} & \textbf{Config Path} & \textbf{Default} & \textbf{Physical Meaning} \\
\hline
\endhead

\texttt{mu} & diffusion.mu & 2.0 & Cell diffusion coefficient \\
\hline

\texttt{nu} & diffusion.nu & 1.0 & Chemical diffusion coefficient \\
\hline

\texttt{a} & reaction.a & 0.0 & Reaction parameter a \\
\hline

\texttt{b} & reaction.b & 1.0 & Reaction parameter b \\
\hline

\texttt{chi} & chemotaxis.chi & 'constant' & Chemotactic sensitivity function \\
\hline

\texttt{dchi} & chemotaxis.dchi & 'zeros' & Derivative of chemotaxis function \\
\hline

\end{longtable}

\subsection{Problem Module Usage Patterns}
\label{subsec:usage_patterns}

\subsubsection{Using OrganOnChip Problem}

\begin{lstlisting}[language=Python, caption=OoC Problem Usage]
# Method 1: Using default configuration
from bionetflux.problems.ooc_problem import create_global_framework
problems, global_disc, constraints, name = create_global_framework()

# Method 2: Using custom geometry
from bionetflux.geometry.domain_geometry import DomainGeometry
geometry = DomainGeometry("custom")
geometry.add_domain(extrema_start=(0,0), extrema_end=(1,0), name="domain1")
problems, global_disc, constraints, name = create_global_framework(geometry=geometry)

# Method 3: Using TOML configuration file
problems, global_disc, constraints, name = create_global_framework(
    config_file="my_ooc_config.toml"
)
\end{lstlisting}

\subsubsection{TOML Configuration Example}

\begin{lstlisting}[language=toml, caption=OoC TOML Configuration]
[problem]
name = "My_OoC_Problem"
neq = 4
problem_type = "ooc"

[time_parameters]
T = 2.0
dt = 0.05

[physical_parameters.viscosity]
nu = 1.5
mu = 2.5
epsilon = 1.2
sigma = 0.8

[physical_parameters.coupling]
chi = 2.0

[initial_conditions]
u = "ooc_gaussian_blob"
v = "ooc_parabolic_flow"

[domain_initial_conditions]
domain_0_u = "constant"
domain_3_v = "ooc_inlet_profile"

[domain_force_functions]
domain_1_omega = "sin_t"
\end{lstlisting}

\subsection{Function Resolution System}
\label{subsec:function_resolution}

\subsubsection{Function Library Hierarchy}

\begin{lstlisting}[language=Python, caption=Function Resolution Process]
# 1. Common functions (in BaseConfigManager)
common_functions = {
    'zeros': lambda s, t=0: np.zeros_like(s),
    'ones': lambda s, t=0: np.ones_like(s),
    'constant': lambda s, t=0: np.ones_like(s),
    'linear': lambda s, t=0: s,
    'quadratic': lambda s, t=0: s**2,
    'sin': lambda s, t=0: np.sin(np.pi * s),
    'cos': lambda s, t=0: np.cos(np.pi * s),
}

# 2. Problem-specific functions (OoC or KS)
# Added by _setup_function_library() in each config manager

# 3. Resolution order:
#    - Check problem-specific library first
#    - Fall back to common library
#    - Raise error if function not found
\end{lstlisting}

\subsection{Validation System}
\label{subsec:validation_system}

\subsubsection{Configuration Validation Rules}

\begin{lstlisting}[language=Python, caption=Validation Rules Example]
# In OoCConfigManager
def _setup_validation_rules(self):
    # Problem type validation
    self.validator.add_rule('problem.problem_type', 
                          {'type': str, 'required': True})
    
    # Physical parameter validation
    self.validator.add_rule('physical_parameters.viscosity.nu', 
                          {'type': float, 'min': 0.0, 'positive': True})
    
    # Time parameter validation
    self.validator.add_rule('time_parameters.T', 
                          {'type': float, 'min': 0.0, 'positive': True})
    
    # Discretization validation
    self.validator.add_rule('discretization.n_elements', 
                          {'type': int, 'min': 1})
\end{lstlisting}

\subsubsection{Problem Type Validation}

\begin{lstlisting}[language=Python, caption=Problem Type Validation]
def load_config(self, config_file: Optional[str] = None):
    config = super().load_config(config_file)
    
    # Validate problem type matches
    if config_file:
        config_problem_type = config.get('problem', {}).get('problem_type', None)
        if config_problem_type and config_problem_type != self.problem_type:
            raise ValueError(
                f"Configuration problem type '{config_problem_type}' does not match "
                f"expected problem type '{self.problem_type}'"
            )
    
    return config
\end{lstlisting}

\subsection{Mathematical Formulation}
\label{subsec:mathematical_formulation}

\subsubsection{OrganOnChip System}

The 4-equation OrganOnChip system implemented in the problems:

\begin{align}
\frac{\partial u}{\partial t} &= \nu \nabla^2 u - \chi \nabla \cdot (u \nabla \phi) + a u + f_u(x,t) \\
\frac{\partial \omega}{\partial t} &= \epsilon \nabla^2 \omega + c \omega + d u + f_\omega(x,t) \\
\frac{\partial v}{\partial t} &= \sigma \nabla^2 v + \lambda(\bar{\omega}) v + f_v(x,t) \\
\frac{\partial \phi}{\partial t} &= \mu \nabla^2 \phi + a \phi + b v + f_\phi(x,t)
\end{align}

where $\bar{\omega} = \frac{1}{|K|} \int_K \omega \, dx$ is the cell-averaged value.

\subsubsection{Keller-Segel System}

The 2-equation Keller-Segel system:

\begin{align}
\frac{\partial u}{\partial t} &= \mu \nabla^2 u - \nabla \cdot (u \chi(\phi) \nabla \phi) + f_u(x,t) \\
\frac{\partial \phi}{\partial t} &= \nu \nabla^2 \phi + u + f_\phi(x,t)
\end{align}

\subsection{Extension Guidelines}
\label{subsec:extension_guidelines}

\subsubsection{Creating New Problem Modules}

\begin{lstlisting}[language=Python, caption=New Problem Module Template]
def create_global_framework(config_file: Optional[str] = None):
    """
    New problem module following BioNetFlux patterns.
    """
    
    # 1. Create specialized config manager
    config_manager = MyNewConfigManager()
    config = config_manager.load_config(config_file)
    
    # 2. Extract parameters and functions
    parameters = extract_parameters_from_config(config)
    initial_conditions = config['initial_conditions']  # Already resolved
    
    # 3. Create geometry (default or provided)
    geometry = build_my_geometry()
    
    # 4. Create problems from geometry
    problems = []
    for domain_id in range(geometry.num_domains()):
        domain_info = geometry.get_domain(domain_id)
        problem = Problem(neq, domain_info.domain_start, domain_info.domain_length,
                         parameters, "my_problem_type", name)
        # Set functions...
        problems.append(problem)
    
    # 5. Setup constraints from geometry
    constraint_manager = setup_constraints_from_geometry(geometry, problems, neq)
    
    # 6. Return framework components
    return problems, global_discretization, constraint_manager, problem_name
\end{lstlisting}

\subsubsection{Configuration Manager Guidelines}

\begin{lstlisting}[language=Python, caption=New Config Manager Template]
class MyNewConfigManager(BaseConfigManager):
    def __init__(self):
        super().__init__("my_type")
    
    def _setup_defaults(self):
        self.defaults = {
            'problem': {'name': 'My_Problem', 'neq': N, 'problem_type': 'my_type'},
            'physical_parameters': { /* your parameters */ },
            'initial_conditions': { /* your variables */ },
            # ...existing code...
        }
    
    def _setup_function_library(self):
        my_functions = {
            'my_special_function': lambda s, t=0: /* implementation */,
            # ...existing code...
        }
        self.function_resolver.register_function_library(my_functions)
    
    def _setup_validation_rules(self):
        # Add your validation rules
        # ...existing code...
\end{lstlisting}

This documentation provides comprehensive coverage of the actual problem modules implemented in the BioNetFlux problems folder, reflecting the current configuration-driven architecture and providing clear guidelines for extension.

% End of problems folder detailed API documentation


\section{Problems Folder Detailed API Reference}
\label{sec:problems_folder_detailed_api}

This section provides comprehensive documentation for the actual problem modules implemented in the BioNetFlux problems folder, based on the current source code analysis.

\subsection{Problems Folder Structure}

The problems folder contains the following implemented modules:

\begin{itemize}
    \item \textbf{custom\_problem\_template.py}: Template for creating new problems
    \item \textbf{ooc\_problem.py}: OrganOnChip grid problem with configuration system
    \item \textbf{ooc\_config\_manager.py}: OrganOnChip-specific configuration manager
    \item \textbf{ks\_config\_manager.py}: Keller-Segel-specific configuration manager
\end{itemize}

\subsection{Configuration-Based Problem Architecture}
\label{subsec:config_architecture}

All current problem modules follow a configuration-driven architecture using TOML files and specialized configuration managers.

\subsubsection{Configuration Manager Pattern}

\begin{lstlisting}[language=Python, caption=Configuration Manager Pattern]
# Each problem type has its own configuration manager
from bionetflux.problems.ooc_config_manager import OoCConfigManager
from bionetflux.problems.ks_config_manager import KSConfigManager

def create_global_framework(config_file: Optional[str] = None):
    # Load configuration with validation
    config_manager = OoCConfigManager()  # or KSConfigManager()
    config = config_manager.load_config(config_file)
    
    # Extract resolved functions from config
    initial_conditions = config['initial_conditions']  # Already callables
    force_functions = config['force_functions']        # Already callables
    # ...existing code...
\end{lstlisting}

\subsection{OrganOnChip Configuration Manager}
\label{subsec:ooc_config_manager}

\subsubsection{OoCConfigManager Class}

\begin{lstlisting}[language=Python, caption=OoC Configuration Manager]
class OoCConfigManager(BaseConfigManager):
    def __init__(self):
        super().__init__("ooc")  # Config manager type
    
    def _setup_defaults(self):
        self.defaults = {
            'problem': {
                'name': 'OoC_Grid_Problem',
                'neq': 4,
                'problem_type': 'ooc'  # Config validation type
            },
            'physical_parameters': {
                'viscosity': {'nu': 1.0, 'mu': 2.0, 'epsilon': 1.0, 'sigma': 1.0},
                'reaction': {'a': 0.0, 'c': 0.0},
                'coupling': {'b': 1.0, 'd': 1.0, 'chi': 1.0}
            },
            'discretization': {
                'n_elements': 20,
                'tau': [0.5, 0.5, 0.5, 0.5]
            },
            'initial_conditions': {
                'u': 'zeros', 'omega': 'zeros', 'v': 'zeros', 'phi': 'zeros'
            },
            'force_functions': {
                'u': 'zeros', 'omega': 'zeros', 'v': 'zeros', 'phi': 'zeros'
            },
            'domain_initial_conditions': {},  # Domain-specific overrides
            'domain_force_functions': {}      # Domain-specific overrides
        }
\end{lstlisting}

\subsubsection{OoC Function Library}

The OoC configuration manager provides specialized functions:

\begin{lstlisting}[language=Python, caption=OoC Function Library]
def _setup_function_library(self):
    ooc_functions = {
        # OrganOnChip specific functions
        'ooc_viscous_profile': lambda s, t=0: s * (1 - s),
        'ooc_parabolic_flow': lambda s, t=0: 4 * s * (1 - s),
        'ooc_inlet_profile': lambda s, t=0: np.where(np.abs(s) < 0.1, 1.0, 0.0),
        
        # Time-dependent functions
        'sin_t': lambda s, t=0: np.sin(t) * np.ones_like(s),
        'cos_t': lambda s, t=0: np.cos(t) * np.ones_like(s),
        'exp_t': lambda s, t=0: np.exp(-t) * np.ones_like(s),
        
        # Combined space-time functions
        'traveling_wave': lambda s, t=0: np.sin(2 * np.pi * (s - t)),
        'diffusive_wave': lambda s, t=0: np.exp(-t) * np.sin(np.pi * s),
    }
\end{lstlisting}

\subsection{Keller-Segel Configuration Manager}
\label{subsec:ks_config_manager}

\subsubsection{KSConfigManager Class}

\begin{lstlisting}[language=Python, caption=KS Configuration Manager]
class KSConfigManager(BaseConfigManager):
    def __init__(self):
        super().__init__("ks")  # Config manager type
    
    def _setup_defaults(self):
        self.defaults = {
            'problem': {
                'name': 'KS_Traveling_Wave_Problem',
                'neq': 2,
                'problem_type': 'ks'
            },
            'physical_parameters': {
                'diffusion': {'mu': 2.0, 'nu': 1.0},
                'reaction': {'a': 0.0, 'b': 1.0},
                'chemotaxis': {
                    'chi': 'constant',    # Function name
                    'dchi': 'zeros'       # Function name
                }
            },
            'discretization': {
                'n_elements': 10,
                'tau': [1.0, 1.0]  # 2 equations
            },
            'initial_conditions': {
                'u': 'zeros',    # Cell density
                'phi': 'zeros'   # Chemical concentration
            }
        }
\end{lstlisting}

\subsubsection{KS Function Library}

\begin{lstlisting}[language=Python, caption=KS Function Library]
def _setup_function_library(self):
    ks_functions = {
        # Keller-Segel specific functions
        'ks_gaussian_blob': lambda s, t=0: np.exp(-(s - 2.0)**2),
        'ks_step_function': lambda s, t=0: np.where(np.abs(s - 2.0) < 0.5, 1.0, 0.0),
        
        # Chemotactic patterns
        'chemotactic_gradient': lambda s, t=0: s * np.exp(-s),
        'attraction_source': lambda s, t=0: np.exp(-((s - 2.0)**2) / 0.5),
    }
\end{lstlisting}

\subsection{Custom Problem Template}
\label{subsec:custom_problem_template}

The \texttt{custom\_problem\_template.py} provides a structured template for creating new problems.

\subsubsection{Template Structure}

\begin{lstlisting}[language=Python, caption=Custom Problem Template Structure]
def create_global_framework():
    """
    Custom problem template following MATLAB TestProblem.m structure.
    Returns: (problems, global_discretization, constraint_manager, problem_name)
    """
    
    # SECTION 1: PHYSICAL PARAMETERS (From MATLAB TestProblem.m)
    neq = 4  # 4-equation OrganOnChip system
    
    # Physical parameters following MATLAB TestProblem.m exactly
    nu, mu, epsilon, sigma = 1.0, 2.0, 1.0, 1.0  # Viscosity
    a, c = 0.0, 0.0                               # Reaction
    b, d, chi = 1.0, 1.0, 1.0                     # Coupling
    parameters = np.array([nu, mu, epsilon, sigma, a, b, c, d, chi])
    
    # SECTION 2: MATHEMATICAL FUNCTIONS
    def lambda_func(x): return np.ones_like(x)    # Constant function
    def dlambda_func(x): return np.zeros_like(x)  # Derivative
    
    # Initial conditions matching MATLAB TestProblem.m
    def initial_u(s, t=0.0): return np.sin(2 * np.pi * s)  # Non-trivial u
    def initial_omega(s, t=0.0): return np.zeros_like(s)    # Zero omega
    def initial_v(s, t=0.0): return np.zeros_like(s)        # Zero v  
    def initial_phi(s, t=0.0): return np.zeros_like(s)      # Zero phi
    
    # Source terms (all zero matching MATLAB)
    def source_u(s, t): return np.zeros_like(s)
    # ... similar for omega, v, phi
    
    # SECTION 3: GEOMETRY SETUP - ⚠️ CUSTOMIZE THIS SECTION ⚠️
    geometry = DomainGeometry("custom_geometry")
    # User adds geometry.add_domain() calls here
    
    # SECTION 4: PROBLEM CREATION (Automated from geometry)
    # SECTION 5: CONSTRAINT SETUP - ⚠️ CUSTOMIZE THIS SECTION ⚠️  
    # SECTION 6-7: GLOBAL DISCRETIZATION AND VALIDATION
\end{lstlisting}

\subsection{OrganOnChip Grid Problem}
\label{subsec:ooc_grid_problem}

The \texttt{ooc\_problem.py} implements a complete OrganOnChip problem with grid geometry.

\subsubsection{Main Framework Method}

\begin{lstlisting}[language=Python, caption=OoC Problem Framework]
def create_global_framework(geometry: Optional[DomainGeometry] = None,
                          config_file: Optional[str] = None):
    """
    OrganOnChip problem with grid geometry and configuration system.
    
    Args:
        geometry: Optional pre-defined geometry. Uses build_default_geometry() if None
        config_file: Optional TOML configuration file
    
    Returns:
        Tuple: (problems, global_discretization, constraint_manager, problem_name)
    """
    
    CONFIG_VALIDATION_TYPE = "ooc"      # For config file validation
    PROBLEM_TYPE = "organ_on_chip"      # For StaticCondensationFactory
    
    # Configuration loading with validation
    config_manager = OoCConfigManager()
    config = config_manager.load_config(config_file)
    
    # Extract configuration sections
    problem_config = config['problem']
    phys_params = config['physical_parameters']
    initial_conditions = config['initial_conditions']  # Already callables
    domain_initial_conditions = config.get('domain_initial_conditions', {})
    domain_force_functions = config.get('domain_force_functions', {})
    
    # Create problems from geometry
    problems = []
    for domain_id in range(geometry.num_domains()):
        domain_info = geometry.get_domain(domain_id)
        problem = Problem(neq, domain_info.domain_start, domain_info.domain_length,
                         parameters, "my_problem_type", name)
        # Set functions...
        problems.append(problem)
    
    # Setup constraints from geometry
    constraint_manager = setup_constraints_from_geometry(geometry, problems, neq)
    
    # Return framework components
    return problems, global_discretization, constraint_manager, problem_name
\end{lstlisting}

\subsubsection{Geometry Integration}

\begin{lstlisting}[language=Python, caption=Geometry Integration]
def build_default_geometry():
    """Build the default OoC grid geometry."""
    return build_grid_geometry()  # From geometry module

def setup_constraints_from_geometry(geometry: DomainGeometry, problems, neq: int):
    """Generate constraint manager from geometry connections."""
    constraint_manager = ConstraintManager()
    
    # Process boundary connections
    boundary_connections = geometry.get_boundary_connections()
    for conn in boundary_connections:
        if conn.is_exterior_boundary():
            for eq_idx in range(neq):
                constraint_manager.add_neumann(
                    equation_index=eq_idx,
                    domain_index=conn.domain1_id,
                    position=conn.parameter1,
                    data_function=lambda t: 0.0  # Homogeneous Neumann
                )
    
    # Process interior connections
    interior_connections = geometry.get_interior_connections()
    for conn in interior_connections:
        for eq_idx in range(neq):
            constraint_manager.add_trace_continuity(
                equation_index=eq_idx,
                domain1_index=conn.domain1_id,
                domain2_index=conn.domain2_id,
                position1=conn.parameter1,
                position2=conn.parameter2
            )
    
    return constraint_manager
\end{lstlisting}

\subsubsection{Domain-Specific Overrides}

\begin{lstlisting}[language=Python, caption=Domain-Specific Configuration]
# Apply domain-specific initial condition overrides
equation_names = ['u', 'omega', 'v', 'phi']

for domain_eq_key, func_name in domain_initial_conditions.items():
    # Parse "domain_<id>_<equation_name>"
    if domain_eq_key.startswith('domain_') and '_' in domain_eq_key[7:]:
        parts = domain_eq_key[7:].split('_', 1)
        domain_id = int(parts[0])
        equation_name = parts[1]
        
        if domain_id < len(problems) and equation_name in equation_names:
            equation_id = equation_names.index(equation_name)
            resolved_func = config_manager.function_resolver.resolve_function(func_name)
            problems[domain_id].set_initial_condition(equation_id, resolved_func)
            print(f"✓ Set initial condition: domain {domain_id}, equation '{equation_name}' -> {func_name}")
\end{lstlisting}

\subsection{Parameter Management}
\label{subsec:parameter_management}

\subsubsection{OrganOnChip Parameter Structure}

\begin{longtable}{|p{3cm}|p{2cm}|p{2cm}|p{6cm}|}
\hline
\textbf{Parameter} & \textbf{Index} & \textbf{Default} & \textbf{Physical Meaning} \\
\hline
\endhead

\texttt{nu} & 0 & 1.0 & Cell diffusion coefficient \\
\hline

\texttt{mu} & 1 & 2.0 & Potential diffusion coefficient \\
\hline

\texttt{epsilon} & 2 & 1.0 & Chemical diffusion coefficient \\
\hline

\texttt{sigma} & 3 & 1.0 & Velocity diffusion coefficient \\
\hline

\texttt{a} & 4 & 0.0 & Growth/decay rate \\
\hline

\texttt{b} & 5 & 1.0 & Velocity-potential coupling \\
\hline

\texttt{c} & 6 & 0.0 & Chemical reaction rate \\
\hline

\texttt{d} & 7 & 1.0 & Cell-chemical coupling \\
\hline

\texttt{chi} & 8 & 1.0 & Chemotactic sensitivity \\
\hline

\end{longtable}

\subsubsection{Keller-Segel Parameter Structure}

\begin{longtable}{|p{3cm}|p{2cm}|p{2cm}|p{6cm}|}
\hline
\textbf{Parameter} & \textbf{Config Path} & \textbf{Default} & \textbf{Physical Meaning} \\
\hline
\endhead

\texttt{mu} & diffusion.mu & 2.0 & Cell diffusion coefficient \\
\hline

\texttt{nu} & diffusion.nu & 1.0 & Chemical diffusion coefficient \\
\hline

\texttt{a} & reaction.a & 0.0 & Reaction parameter a \\
\hline

\texttt{b} & reaction.b & 1.0 & Reaction parameter b \\
\hline

\texttt{chi} & chemotaxis.chi & 'constant' & Chemotactic sensitivity function \\
\hline

\texttt{dchi} & chemotaxis.dchi & 'zeros' & Derivative of chemotaxis function \\
\hline

\end{longtable}

\subsection{Problem Module Usage Patterns}
\label{subsec:usage_patterns}

\subsubsection{Using OrganOnChip Problem}

\begin{lstlisting}[language=Python, caption=OoC Problem Usage]
# Method 1: Using default configuration
from bionetflux.problems.ooc_problem import create_global_framework
problems, global_disc, constraints, name = create_global_framework()

# Method 2: Using custom geometry
from bionetflux.geometry.domain_geometry import DomainGeometry
geometry = DomainGeometry("custom")
geometry.add_domain(extrema_start=(0,0), extrema_end=(1,0), name="domain1")
problems, global_disc, constraints, name = create_global_framework(geometry=geometry)

# Method 3: Using TOML configuration file
problems, global_disc, constraints, name = create_global_framework(
    config_file="my_ooc_config.toml"
)
\end{lstlisting}

\subsubsection{TOML Configuration Example}

\begin{lstlisting}[language=toml, caption=OoC TOML Configuration]
[problem]
name = "My_OoC_Problem"
neq = 4
problem_type = "ooc"

[time_parameters]
T = 2.0
dt = 0.05

[physical_parameters.viscosity]
nu = 1.5
mu = 2.5
epsilon = 1.2
sigma = 0.8

[physical_parameters.coupling]
chi = 2.0

[initial_conditions]
u = "ooc_gaussian_blob"
v = "ooc_parabolic_flow"

[domain_initial_conditions]
domain_0_u = "constant"
domain_3_v = "ooc_inlet_profile"

[domain_force_functions]
domain_1_omega = "sin_t"
\end{lstlisting}

\subsection{Function Resolution System}
\label{subsec:function_resolution}

\subsubsection{Function Library Hierarchy}

\begin{lstlisting}[language=Python, caption=Function Resolution Process]
# 1. Common functions (in BaseConfigManager)
common_functions = {
    'zeros': lambda s, t=0: np.zeros_like(s),
    'ones': lambda s, t=0: np.ones_like(s),
    'constant': lambda s, t=0: np.ones_like(s),
    'linear': lambda s, t=0: s,
    'quadratic': lambda s, t=0: s**2,
    'sin': lambda s, t=0: np.sin(np.pi * s),
    'cos': lambda s, t=0: np.cos(np.pi * s),
}

# 2. Problem-specific functions (OoC or KS)
# Added by _setup_function_library() in each config manager

# 3. Resolution order:
#    - Check problem-specific library first
#    - Fall back to common library
#    - Raise error if function not found
\end{lstlisting}

\subsection{Validation System}
\label{subsec:validation_system}

\subsubsection{Configuration Validation Rules}

\begin{lstlisting}[language=Python, caption=Validation Rules Example]
# In OoCConfigManager
def _setup_validation_rules(self):
    # Problem type validation
    self.validator.add_rule('problem.problem_type', 
                          {'type': str, 'required': True})
    
    # Physical parameter validation
    self.validator.add_rule('physical_parameters.viscosity.nu', 
                          {'type': float, 'min': 0.0, 'positive': True})
    
    # Time parameter validation
    self.validator.add_rule('time_parameters.T', 
                          {'type': float, 'min': 0.0, 'positive': True})
    
    # Discretization validation
    self.validator.add_rule('discretization.n_elements', 
                          {'type': int, 'min': 1})
\end{lstlisting}

\subsubsection{Problem Type Validation}

\begin{lstlisting}[language=Python, caption=Problem Type Validation]
def load_config(self, config_file: Optional[str] = None):
    config = super().load_config(config_file)
    
    # Validate problem type matches
    if config_file:
        config_problem_type = config.get('problem', {}).get('problem_type', None)
        if config_problem_type and config_problem_type != self.problem_type:
            raise ValueError(
                f"Configuration problem type '{config_problem_type}' does not match "
                f"expected problem type '{self.problem_type}'"
            )
    
    return config
\end{lstlisting}

\subsection{Mathematical Formulation}
\label{subsec:mathematical_formulation}

\subsubsection{OrganOnChip System}

The 4-equation OrganOnChip system implemented in the problems:

\begin{align}
\frac{\partial u}{\partial t} &= \nu \nabla^2 u - \chi \nabla \cdot (u \nabla \phi) + a u + f_u(x,t) \\
\frac{\partial \omega}{\partial t} &= \epsilon \nabla^2 \omega + c \omega + d u + f_\omega(x,t) \\
\frac{\partial v}{\partial t} &= \sigma \nabla^2 v + \lambda(\bar{\omega}) v + f_v(x,t) \\
\frac{\partial \phi}{\partial t} &= \mu \nabla^2 \phi + a \phi + b v + f_\phi(x,t)
\end{align}

where $\bar{\omega} = \frac{1}{|K|} \int_K \omega \, dx$ is the cell-averaged value.

\subsubsection{Keller-Segel System}

The 2-equation Keller-Segel system:

\begin{align}
\frac{\partial u}{\partial t} &= \mu \nabla^2 u - \nabla \cdot (u \chi(\phi) \nabla \phi) + f_u(x,t) \\
\frac{\partial \phi}{\partial t} &= \nu \nabla^2 \phi + u + f_\phi(x,t)
\end{align}

\subsection{Extension Guidelines}
\label{subsec:extension_guidelines}

\subsubsection{Creating New Problem Modules}

\begin{lstlisting}[language=Python, caption=New Problem Module Template]
def create_global_framework(config_file: Optional[str] = None):
    """
    New problem module following BioNetFlux patterns.
    """
    
    # 1. Create specialized config manager
    config_manager = MyNewConfigManager()
    config = config_manager.load_config(config_file)
    
    # 2. Extract parameters and functions
    parameters = extract_parameters_from_config(config)
    initial_conditions = config['initial_conditions']  # Already resolved
    
    # 3. Create geometry (default or provided)
    geometry = build_my_geometry()
    
    # 4. Create problems from geometry
    problems = []
    for domain_id in range(geometry.num_domains()):
        domain_info = geometry.get_domain(domain_id)
        problem = Problem(neq, domain_info.domain_start, domain_info.domain_length,
                         parameters, "my_problem_type", name)
        # Set functions...
        problems.append(problem)
    
    # 5. Setup constraints from geometry
    constraint_manager = setup_constraints_from_geometry(geometry, problems, neq)
    
    # 6. Return framework components
    return problems, global_discretization, constraint_manager, problem_name
\end{lstlisting}

\subsubsection{Configuration Manager Guidelines}

\begin{lstlisting}[language=Python, caption=New Config Manager Template]
class MyNewConfigManager(BaseConfigManager):
    def __init__(self):
        super().__init__("my_type")
    
    def _setup_defaults(self):
        self.defaults = {
            'problem': {'name': 'My_Problem', 'neq': N, 'problem_type': 'my_type'},
            'physical_parameters': { /* your parameters */ },
            'initial_conditions': { /* your variables */ },
            # ...existing code...
        }
    
    def _setup_function_library(self):
        my_functions = {
            'my_special_function': lambda s, t=0: /* implementation */,
            # ...existing code...
        }
        self.function_resolver.register_function_library(my_functions)
    
    def _setup_validation_rules(self):
        # Add your validation rules
        # ...existing code...
\end{lstlisting}

This documentation provides comprehensive coverage of the actual problem modules implemented in the BioNetFlux problems folder, reflecting the current configuration-driven architecture and providing clear guidelines for extension.

% End of problems folder detailed API documentation


\section{Problems Folder Detailed API Reference}
\label{sec:problems_folder_detailed_api}

This section provides comprehensive documentation for the actual problem modules implemented in the BioNetFlux problems folder, based on the current source code analysis.

\subsection{Problems Folder Structure}

The problems folder contains the following implemented modules:

\begin{itemize}
    \item \textbf{custom\_problem\_template.py}: Template for creating new problems
    \item \textbf{ooc\_problem.py}: OrganOnChip grid problem with configuration system
    \item \textbf{ooc\_config\_manager.py}: OrganOnChip-specific configuration manager
    \item \textbf{ks\_config\_manager.py}: Keller-Segel-specific configuration manager
\end{itemize}

\subsection{Configuration-Based Problem Architecture}
\label{subsec:config_architecture}

All current problem modules follow a configuration-driven architecture using TOML files and specialized configuration managers.

\subsubsection{Configuration Manager Pattern}

\begin{lstlisting}[language=Python, caption=Configuration Manager Pattern]
# Each problem type has its own configuration manager
from bionetflux.problems.ooc_config_manager import OoCConfigManager
from bionetflux.problems.ks_config_manager import KSConfigManager

def create_global_framework(config_file: Optional[str] = None):
    # Load configuration with validation
    config_manager = OoCConfigManager()  # or KSConfigManager()
    config = config_manager.load_config(config_file)
    
    # Extract resolved functions from config
    initial_conditions = config['initial_conditions']  # Already callables
    force_functions = config['force_functions']        # Already callables
    # ...existing code...
\end{lstlisting}

\subsection{OrganOnChip Configuration Manager}
\label{subsec:ooc_config_manager}

\subsubsection{OoCConfigManager Class}

\begin{lstlisting}[language=Python, caption=OoC Configuration Manager]
class OoCConfigManager(BaseConfigManager):
    def __init__(self):
        super().__init__("ooc")  # Config manager type
    
    def _setup_defaults(self):
        self.defaults = {
            'problem': {
                'name': 'OoC_Grid_Problem',
                'neq': 4,
                'problem_type': 'ooc'  # Config validation type
            },
            'physical_parameters': {
                'viscosity': {'nu': 1.0, 'mu': 2.0, 'epsilon': 1.0, 'sigma': 1.0},
                'reaction': {'a': 0.0, 'c': 0.0},
                'coupling': {'b': 1.0, 'd': 1.0, 'chi': 1.0}
            },
            'discretization': {
                'n_elements': 20,
                'tau': [0.5, 0.5, 0.5, 0.5]
            },
            'initial_conditions': {
                'u': 'zeros', 'omega': 'zeros', 'v': 'zeros', 'phi': 'zeros'
            },
            'force_functions': {
                'u': 'zeros', 'omega': 'zeros', 'v': 'zeros', 'phi': 'zeros'
            },
            'domain_initial_conditions': {},  # Domain-specific overrides
            'domain_force_functions': {}      # Domain-specific overrides
        }
\end{lstlisting}

\subsubsection{OoC Function Library}

The OoC configuration manager provides specialized functions:

\begin{lstlisting}[language=Python, caption=OoC Function Library]
def _setup_function_library(self):
    ooc_functions = {
        # OrganOnChip specific functions
        'ooc_viscous_profile': lambda s, t=0: s * (1 - s),
        'ooc_parabolic_flow': lambda s, t=0: 4 * s * (1 - s),
        'ooc_inlet_profile': lambda s, t=0: np.where(np.abs(s) < 0.1, 1.0, 0.0),
        
        # Time-dependent functions
        'sin_t': lambda s, t=0: np.sin(t) * np.ones_like(s),
        'cos_t': lambda s, t=0: np.cos(t) * np.ones_like(s),
        'exp_t': lambda s, t=0: np.exp(-t) * np.ones_like(s),
        
        # Combined space-time functions
        'traveling_wave': lambda s, t=0: np.sin(2 * np.pi * (s - t)),
        'diffusive_wave': lambda s, t=0: np.exp(-t) * np.sin(np.pi * s),
    }
\end{lstlisting}

\subsection{Keller-Segel Configuration Manager}
\label{subsec:ks_config_manager}

\subsubsection{KSConfigManager Class}

\begin{lstlisting}[language=Python, caption=KS Configuration Manager]
class KSConfigManager(BaseConfigManager):
    def __init__(self):
        super().__init__("ks")  # Config manager type
    
    def _setup_defaults(self):
        self.defaults = {
            'problem': {
                'name': 'KS_Traveling_Wave_Problem',
                'neq': 2,
                'problem_type': 'ks'
            },
            'physical_parameters': {
                'diffusion': {'mu': 2.0, 'nu': 1.0},
                'reaction': {'a': 0.0, 'b': 1.0},
                'chemotaxis': {
                    'chi': 'constant',    # Function name
                    'dchi': 'zeros'       # Function name
                }
            },
            'discretization': {
                'n_elements': 10,
                'tau': [1.0, 1.0]  # 2 equations
            },
            'initial_conditions': {
                'u': 'zeros',    # Cell density
                'phi': 'zeros'   # Chemical concentration
            }
        }
\end{lstlisting}

\subsubsection{KS Function Library}

\begin{lstlisting}[language=Python, caption=KS Function Library]
def _setup_function_library(self):
    ks_functions = {
        # Keller-Segel specific functions
        'ks_gaussian_blob': lambda s, t=0: np.exp(-(s - 2.0)**2),
        'ks_step_function': lambda s, t=0: np.where(np.abs(s - 2.0) < 0.5, 1.0, 0.0),
        
        # Chemotactic patterns
        'chemotactic_gradient': lambda s, t=0: s * np.exp(-s),
        'attraction_source': lambda s, t=0: np.exp(-((s - 2.0)**2) / 0.5),
    }
\end{lstlisting}

\subsection{Custom Problem Template}
\label{subsec:custom_problem_template}

The \texttt{custom\_problem\_template.py} provides a structured template for creating new problems.

\subsubsection{Template Structure}

\begin{lstlisting}[language=Python, caption=Custom Problem Template Structure]
def create_global_framework():
    """
    Custom problem template following MATLAB TestProblem.m structure.
    Returns: (problems, global_discretization, constraint_manager, problem_name)
    """
    
    # SECTION 1: PHYSICAL PARAMETERS (From MATLAB TestProblem.m)
    neq = 4  # 4-equation OrganOnChip system
    
    # Physical parameters following MATLAB TestProblem.m exactly
    nu, mu, epsilon, sigma = 1.0, 2.0, 1.0, 1.0  # Viscosity
    a, c = 0.0, 0.0                               # Reaction
    b, d, chi = 1.0, 1.0, 1.0                     # Coupling
    parameters = np.array([nu, mu, epsilon, sigma, a, b, c, d, chi])
    
    # SECTION 2: MATHEMATICAL FUNCTIONS
    def lambda_func(x): return np.ones_like(x)    # Constant function
    def dlambda_func(x): return np.zeros_like(x)  # Derivative
    
    # Initial conditions matching MATLAB TestProblem.m
    def initial_u(s, t=0.0): return np.sin(2 * np.pi * s)  # Non-trivial u
    def initial_omega(s, t=0.0): return np.zeros_like(s)    # Zero omega
    def initial_v(s, t=0.0): return np.zeros_like(s)        # Zero v  
    def initial_phi(s, t=0.0): return np.zeros_like(s)      # Zero phi
    
    # Source terms (all zero matching MATLAB)
    def source_u(s, t): return np.zeros_like(s)
    # ... similar for omega, v, phi
    
    # SECTION 3: GEOMETRY SETUP - ⚠️ CUSTOMIZE THIS SECTION ⚠️
    geometry = DomainGeometry("custom_geometry")
    # User adds geometry.add_domain() calls here
    
    # SECTION 4: PROBLEM CREATION (Automated from geometry)
    # SECTION 5: CONSTRAINT SETUP - ⚠️ CUSTOMIZE THIS SECTION ⚠️  
    # SECTION 6-7: GLOBAL DISCRETIZATION AND VALIDATION
\end{lstlisting}

\subsection{OrganOnChip Grid Problem}
\label{subsec:ooc_grid_problem}

The \texttt{ooc\_problem.py} implements a complete OrganOnChip problem with grid geometry.

\subsubsection{Main Framework Method}

\begin{lstlisting}[language=Python, caption=OoC Problem Framework]
def create_global_framework(geometry: Optional[DomainGeometry] = None,
                          config_file: Optional[str] = None):
    """
    OrganOnChip problem with grid geometry and configuration system.
    
    Args:
        geometry: Optional pre-defined geometry. Uses build_default_geometry() if None
        config_file: Optional TOML configuration file
    
    Returns:
        Tuple: (problems, global_discretization, constraint_manager, problem_name)
    """
    
    CONFIG_VALIDATION_TYPE = "ooc"      # For config file validation
    PROBLEM_TYPE = "organ_on_chip"      # For StaticCondensationFactory
    
    # Configuration loading with validation
    config_manager = OoCConfigManager()
    config = config_manager.load_config(config_file)
    
    # Extract configuration sections
    problem_config = config['problem']
    phys_params = config['physical_parameters']
    initial_conditions = config['initial_conditions']  # Already callables
    domain_initial_conditions = config.get('domain_initial_conditions', {})
    domain_force_functions = config.get('domain_force_functions', {})
    
    # Create problems from geometry
    problems = []
    for domain_id in range(geometry.num_domains()):
        domain_info = geometry.get_domain(domain_id)
        problem = Problem(neq, domain_info.domain_start, domain_info.domain_length,
                         parameters, "my_problem_type", name)
        # Set functions...
        problems.append(problem)
    
    # Setup constraints from geometry
    constraint_manager = setup_constraints_from_geometry(geometry, problems, neq)
    
    # Return framework components
    return problems, global_discretization, constraint_manager, problem_name
\end{lstlisting}

\subsubsection{Geometry Integration}

\begin{lstlisting}[language=Python, caption=Geometry Integration]
def build_default_geometry():
    """Build the default OoC grid geometry."""
    return build_grid_geometry()  # From geometry module

def setup_constraints_from_geometry(geometry: DomainGeometry, problems, neq: int):
    """Generate constraint manager from geometry connections."""
    constraint_manager = ConstraintManager()
    
    # Process boundary connections
    boundary_connections = geometry.get_boundary_connections()
    for conn in boundary_connections:
        if conn.is_exterior_boundary():
            for eq_idx in range(neq):
                constraint_manager.add_neumann(
                    equation_index=eq_idx,
                    domain_index=conn.domain1_id,
                    position=conn.parameter1,
                    data_function=lambda t: 0.0  # Homogeneous Neumann
                )
    
    # Process interior connections
    interior_connections = geometry.get_interior_connections()
    for conn in interior_connections:
        for eq_idx in range(neq):
            constraint_manager.add_trace_continuity(
                equation_index=eq_idx,
                domain1_index=conn.domain1_id,
                domain2_index=conn.domain2_id,
                position1=conn.parameter1,
                position2=conn.parameter2
            )
    
    return constraint_manager
\end{lstlisting}

\subsubsection{Domain-Specific Overrides}

\begin{lstlisting}[language=Python, caption=Domain-Specific Configuration]
# Apply domain-specific initial condition overrides
equation_names = ['u', 'omega', 'v', 'phi']

for domain_eq_key, func_name in domain_initial_conditions.items():
    # Parse "domain_<id>_<equation_name>"
    if domain_eq_key.startswith('domain_') and '_' in domain_eq_key[7:]:
        parts = domain_eq_key[7:].split('_', 1)
        domain_id = int(parts[0])
        equation_name = parts[1]
        
        if domain_id < len(problems) and equation_name in equation_names:
            equation_id = equation_names.index(equation_name)
            resolved_func = config_manager.function_resolver.resolve_function(func_name)
            problems[domain_id].set_initial_condition(equation_id, resolved_func)
            print(f"✓ Set initial condition: domain {domain_id}, equation '{equation_name}' -> {func_name}")
\end{lstlisting}

\subsection{Parameter Management}
\label{subsec:parameter_management}

\subsubsection{OrganOnChip Parameter Structure}

\begin{longtable}{|p{3cm}|p{2cm}|p{2cm}|p{6cm}|}
\hline
\textbf{Parameter} & \textbf{Index} & \textbf{Default} & \textbf{Physical Meaning} \\
\hline
\endhead

\texttt{nu} & 0 & 1.0 & Cell diffusion coefficient \\
\hline

\texttt{mu} & 1 & 2.0 & Potential diffusion coefficient \\
\hline

\texttt{epsilon} & 2 & 1.0 & Chemical diffusion coefficient \\
\hline

\texttt{sigma} & 3 & 1.0 & Velocity diffusion coefficient \\
\hline

\texttt{a} & 4 & 0.0 & Growth/decay rate \\
\hline

\texttt{b} & 5 & 1.0 & Velocity-potential coupling \\
\hline

\texttt{c} & 6 & 0.0 & Chemical reaction rate \\
\hline

\texttt{d} & 7 & 1.0 & Cell-chemical coupling \\
\hline

\texttt{chi} & 8 & 1.0 & Chemotactic sensitivity \\
\hline

\end{longtable}

\subsubsection{Keller-Segel Parameter Structure}

\begin{longtable}{|p{3cm}|p{2cm}|p{2cm}|p{6cm}|}
\hline
\textbf{Parameter} & \textbf{Config Path} & \textbf{Default} & \textbf{Physical Meaning} \\
\hline
\endhead

\texttt{mu} & diffusion.mu & 2.0 & Cell diffusion coefficient \\
\hline

\texttt{nu} & diffusion.nu & 1.0 & Chemical diffusion coefficient \\
\hline

\texttt{a} & reaction.a & 0.0 & Reaction parameter a \\
\hline

\texttt{b} & reaction.b & 1.0 & Reaction parameter b \\
\hline

\texttt{chi} & chemotaxis.chi & 'constant' & Chemotactic sensitivity function \\
\hline

\texttt{dchi} & chemotaxis.dchi & 'zeros' & Derivative of chemotaxis function \\
\hline

\end{longtable}

\subsection{Problem Module Usage Patterns}
\label{subsec:usage_patterns}

\subsubsection{Using OrganOnChip Problem}

\begin{lstlisting}[language=Python, caption=OoC Problem Usage]
# Method 1: Using default configuration
from bionetflux.problems.ooc_problem import create_global_framework
problems, global_disc, constraints, name = create_global_framework()

# Method 2: Using custom geometry
from bionetflux.geometry.domain_geometry import DomainGeometry
geometry = DomainGeometry("custom")
geometry.add_domain(extrema_start=(0,0), extrema_end=(1,0), name="domain1")
problems, global_disc, constraints, name = create_global_framework(geometry=geometry)

# Method 3: Using TOML configuration file
problems, global_disc, constraints, name = create_global_framework(
    config_file="my_ooc_config.toml"
)
\end{lstlisting}

\subsubsection{TOML Configuration Example}

\begin{lstlisting}[language=toml, caption=OoC TOML Configuration]
[problem]
name = "My_OoC_Problem"
neq = 4
problem_type = "ooc"

[time_parameters]
T = 2.0
dt = 0.05

[physical_parameters.viscosity]
nu = 1.5
mu = 2.5
epsilon = 1.2
sigma = 0.8

[physical_parameters.coupling]
chi = 2.0

[initial_conditions]
u = "ooc_gaussian_blob"
v = "ooc_parabolic_flow"

[domain_initial_conditions]
domain_0_u = "constant"
domain_3_v = "ooc_inlet_profile"

[domain_force_functions]
domain_1_omega = "sin_t"
\end{lstlisting}

\subsection{Function Resolution System}
\label{subsec:function_resolution}

\subsubsection{Function Library Hierarchy}

\begin{lstlisting}[language=Python, caption=Function Resolution Process]
# 1. Common functions (in BaseConfigManager)
common_functions = {
    'zeros': lambda s, t=0: np.zeros_like(s),
    'ones': lambda s, t=0: np.ones_like(s),
    'constant': lambda s, t=0: np.ones_like(s),
    'linear': lambda s, t=0: s,
    'quadratic': lambda s, t=0: s**2,
    'sin': lambda s, t=0: np.sin(np.pi * s),
    'cos': lambda s, t=0: np.cos(np.pi * s),
}

# 2. Problem-specific functions (OoC or KS)
# Added by _setup_function_library() in each config manager

# 3. Resolution order:
#    - Check problem-specific library first
#    - Fall back to common library
#    - Raise error if function not found
\end{lstlisting}

\subsection{Validation System}
\label{subsec:validation_system}

\subsubsection{Configuration Validation Rules}

\begin{lstlisting}[language=Python, caption=Validation Rules Example]
# In OoCConfigManager
def _setup_validation_rules(self):
    # Problem type validation
    self.validator.add_rule('problem.problem_type', 
                          {'type': str, 'required': True})
    
    # Physical parameter validation
    self.validator.add_rule('physical_parameters.viscosity.nu', 
                          {'type': float, 'min': 0.0, 'positive': True})
    
    # Time parameter validation
    self.validator.add_rule('time_parameters.T', 
                          {'type': float, 'min': 0.0, 'positive': True})
    
    # Discretization validation
    self.validator.add_rule('discretization.n_elements', 
                          {'type': int, 'min': 1})
\end{lstlisting}

\subsubsection{Problem Type Validation}

\begin{lstlisting}[language=Python, caption=Problem Type Validation]
def load_config(self, config_file: Optional[str] = None):
    config = super().load_config(config_file)
    
    # Validate problem type matches
    if config_file:
        config_problem_type = config.get('problem', {}).get('problem_type', None)
        if config_problem_type and config_problem_type != self.problem_type:
            raise ValueError(
                f"Configuration problem type '{config_problem_type}' does not match "
                f"expected problem type '{self.problem_type}'"
            )
    
    return config
\end{lstlisting}

\subsection{Mathematical Formulation}
\label{subsec:mathematical_formulation}

\subsubsection{OrganOnChip System}

The 4-equation OrganOnChip system implemented in the problems:

\begin{align}
\frac{\partial u}{\partial t} &= \nu \nabla^2 u - \chi \nabla \cdot (u \nabla \phi) + a u + f_u(x,t) \\
\frac{\partial \omega}{\partial t} &= \epsilon \nabla^2 \omega + c \omega + d u + f_\omega(x,t) \\
\frac{\partial v}{\partial t} &= \sigma \nabla^2 v + \lambda(\bar{\omega}) v + f_v(x,t) \\
\frac{\partial \phi}{\partial t} &= \mu \nabla^2 \phi + a \phi + b v + f_\phi(x,t)
\end{align}

where $\bar{\omega} = \frac{1}{|K|} \int_K \omega \, dx$ is the cell-averaged value.

\subsubsection{Keller-Segel System}

The 2-equation Keller-Segel system:

\begin{align}
\frac{\partial u}{\partial t} &= \mu \nabla^2 u - \nabla \cdot (u \chi(\phi) \nabla \phi) + f_u(x,t) \\
\frac{\partial \phi}{\partial t} &= \nu \nabla^2 \phi + u + f_\phi(x,t)
\end{align}

\subsection{Extension Guidelines}
\label{subsec:extension_guidelines}

\subsubsection{Creating New Problem Modules}

\begin{lstlisting}[language=Python, caption=New Problem Module Template]
def create_global_framework(config_file: Optional[str] = None):
    """
    New problem module following BioNetFlux patterns.
    """
    
    # 1. Create specialized config manager
    config_manager = MyNewConfigManager()
    config = config_manager.load_config(config_file)
    
    # 2. Extract parameters and functions
    parameters = extract_parameters_from_config(config)
    initial_conditions = config['initial_conditions']  # Already resolved
    
    # 3. Create geometry (default or provided)
    geometry = build_my_geometry()
    
    # 4. Create problems from geometry
    problems = []
    for domain_id in range(geometry.num_domains()):
        domain_info = geometry.get_domain(domain_id)
        problem = Problem(neq, domain_info.domain_start, domain_info.domain_length,
                         parameters, "my_problem_type", name)
        # Set functions...
        problems.append(problem)
    
    # 5. Setup constraints from geometry
    constraint_manager = setup_constraints_from_geometry(geometry, problems, neq)
    
    # 6. Return framework components
    return problems, global_discretization, constraint_manager, problem_name
\end{lstlisting}

\subsubsection{Configuration Manager Guidelines}

\begin{lstlisting}[language=Python, caption=New Config Manager Template]
class MyNewConfigManager(BaseConfigManager):
    def __init__(self):
        super().__init__("my_type")
    
    def _setup_defaults(self):
        self.defaults = {
            'problem': {'name': 'My_Problem', 'neq': N, 'problem_type': 'my_type'},
            'physical_parameters': { /* your parameters */ },
            'initial_conditions': { /* your variables */ },
            # ...existing code...
        }
    
    def _setup_function_library(self):
        my_functions = {
            'my_special_function': lambda s, t=0: /* implementation */,
            # ...existing code...
        }
        self.function_resolver.register_function_library(my_functions)
    
    def _setup_validation_rules(self):
        # Add your validation rules
        # ...existing code...
\end{lstlisting}

This documentation provides comprehensive coverage of the actual problem modules implemented in the BioNetFlux problems folder, reflecting the current configuration-driven architecture and providing clear guidelines for extension.

% End of problems folder detailed API documentation


\section{Problems Folder Detailed API Reference}
\label{sec:problems_folder_detailed_api}

This section provides comprehensive documentation for the actual problem modules implemented in the BioNetFlux problems folder, based on the current source code analysis.

\subsection{Problems Folder Structure}

The problems folder contains the following implemented modules:

\begin{itemize}
    \item \textbf{custom\_problem\_template.py}: Template for creating new problems
    \item \textbf{ooc\_problem.py}: OrganOnChip grid problem with configuration system
    \item \textbf{ooc\_config\_manager.py}: OrganOnChip-specific configuration manager
    \item \textbf{ks\_config\_manager.py}: Keller-Segel-specific configuration manager
\end{itemize}

\subsection{Configuration-Based Problem Architecture}
\label{subsec:config_architecture}

All current problem modules follow a configuration-driven architecture using TOML files and specialized configuration managers.

\subsubsection{Configuration Manager Pattern}

\begin{lstlisting}[language=Python, caption=Configuration Manager Pattern]
# Each problem type has its own configuration manager
from bionetflux.problems.ooc_config_manager import OoCConfigManager
from bionetflux.problems.ks_config_manager import KSConfigManager

def create_global_framework(config_file: Optional[str] = None):
    # Load configuration with validation
    config_manager = OoCConfigManager()  # or KSConfigManager()
    config = config_manager.load_config(config_file)
    
    # Extract resolved functions from config
    initial_conditions = config['initial_conditions']  # Already callables
    force_functions = config['force_functions']        # Already callables
    # ...existing code...
\end{lstlisting}

\subsection{OrganOnChip Configuration Manager}
\label{subsec:ooc_config_manager}

\subsubsection{OoCConfigManager Class}

\begin{lstlisting}[language=Python, caption=OoC Configuration Manager]
class OoCConfigManager(BaseConfigManager):
    def __init__(self):
        super().__init__("ooc")  # Config manager type
    
    def _setup_defaults(self):
        self.defaults = {
            'problem': {
                'name': 'OoC_Grid_Problem',
                'neq': 4,
                'problem_type': 'ooc'  # Config validation type
            },
            'physical_parameters': {
                'viscosity': {'nu': 1.0, 'mu': 2.0, 'epsilon': 1.0, 'sigma': 1.0},
                'reaction': {'a': 0.0, 'c': 0.0},
                'coupling': {'b': 1.0, 'd': 1.0, 'chi': 1.0}
            },
            'discretization': {
                'n_elements': 20,
                'tau': [0.5, 0.5, 0.5, 0.5]
            },
            'initial_conditions': {
                'u': 'zeros', 'omega': 'zeros', 'v': 'zeros', 'phi': 'zeros'
            },
            'force_functions': {
                'u': 'zeros', 'omega': 'zeros', 'v': 'zeros', 'phi': 'zeros'
            },
            'domain_initial_conditions': {},  # Domain-specific overrides
            'domain_force_functions': {}      # Domain-specific overrides
        }
\end{lstlisting}

\subsubsection{OoC Function Library}

The OoC configuration manager provides specialized functions:

\begin{lstlisting}[language=Python, caption=OoC Function Library]
def _setup_function_library(self):
    ooc_functions = {
        # OrganOnChip specific functions
        'ooc_viscous_profile': lambda s, t=0: s * (1 - s),
        'ooc_parabolic_flow': lambda s, t=0: 4 * s * (1 - s),
        'ooc_inlet_profile': lambda s, t=0: np.where(np.abs(s) < 0.1, 1.0, 0.0),
        
        # Time-dependent functions
        'sin_t': lambda s, t=0: np.sin(t) * np.ones_like(s),
        'cos_t': lambda s, t=0: np.cos(t) * np.ones_like(s),
        'exp_t': lambda s, t=0: np.exp(-t) * np.ones_like(s),
        
        # Combined space-time functions
        'traveling_wave': lambda s, t=0: np.sin(2 * np.pi * (s - t)),
        'diffusive_wave': lambda s, t=0: np.exp(-t) * np.sin(np.pi * s),
    }
\end{lstlisting}

\subsection{Keller-Segel Configuration Manager}
\label{subsec:ks_config_manager}

\subsubsection{KSConfigManager Class}

\begin{lstlisting}[language=Python, caption=KS Configuration Manager]
class KSConfigManager(BaseConfigManager):
    def __init__(self):
        super().__init__("ks")  # Config manager type
    
    def _setup_defaults(self):
        self.defaults = {
            'problem': {
                'name': 'KS_Traveling_Wave_Problem',
                'neq': 2,
                'problem_type': 'ks'
            },
            'physical_parameters': {
                'diffusion': {'mu': 2.0, 'nu': 1.0},
                'reaction': {'a': 0.0, 'b': 1.0},
                'chemotaxis': {
                    'chi': 'constant',    # Function name
                    'dchi': 'zeros'       # Function name
                }
            },
            'discretization': {
                'n_elements': 10,
                'tau': [1.0, 1.0]  # 2 equations
            },
            'initial_conditions': {
                'u': 'zeros',    # Cell density
                'phi': 'zeros'   # Chemical concentration
            }
        }
\end{lstlisting}

\subsubsection{KS Function Library}

\begin{lstlisting}[language=Python, caption=KS Function Library]
def _setup_function_library(self):
    ks_functions = {
        # Keller-Segel specific functions
        'ks_gaussian_blob': lambda s, t=0: np.exp(-(s - 2.0)**2),
        'ks_step_function': lambda s, t=0: np.where(np.abs(s - 2.0) < 0.5, 1.0, 0.0),
        
        # Chemotactic patterns
        'chemotactic_gradient': lambda s, t=0: s * np.exp(-s),
        'attraction_source': lambda s, t=0: np.exp(-((s - 2.0)**2) / 0.5),
    }
\end{lstlisting}

\subsection{Custom Problem Template}
\label{subsec:custom_problem_template}

The \texttt{custom\_problem\_template.py} provides a structured template for creating new problems.

\subsubsection{Template Structure}

\begin{lstlisting}[language=Python, caption=Custom Problem Template Structure]
def create_global_framework():
    """
    Custom problem template following MATLAB TestProblem.m structure.
    Returns: (problems, global_discretization, constraint_manager, problem_name)
    """
    
    # SECTION 1: PHYSICAL PARAMETERS (From MATLAB TestProblem.m)
    neq = 4  # 4-equation OrganOnChip system
    
    # Physical parameters following MATLAB TestProblem.m exactly
    nu, mu, epsilon, sigma = 1.0, 2.0, 1.0, 1.0  # Viscosity
    a, c = 0.0, 0.0                               # Reaction
    b, d, chi = 1.0, 1.0, 1.0                     # Coupling
    parameters = np.array([nu, mu, epsilon, sigma, a, b, c, d, chi])
    
    # SECTION 2: MATHEMATICAL FUNCTIONS
    def lambda_func(x): return np.ones_like(x)    # Constant function
    def dlambda_func(x): return np.zeros_like(x)  # Derivative
    
    # Initial conditions matching MATLAB TestProblem.m
    def initial_u(s, t=0.0): return np.sin(2 * np.pi * s)  # Non-trivial u
    def initial_omega(s, t=0.0): return np.zeros_like(s)    # Zero omega
    def initial_v(s, t=0.0): return np.zeros_like(s)        # Zero v  
    def initial_phi(s, t=0.0): return np.zeros_like(s)      # Zero phi
    
    # Source terms (all zero matching MATLAB)
    def source_u(s, t): return np.zeros_like(s)
    # ... similar for omega, v, phi
    
    # SECTION 3: GEOMETRY SETUP - ⚠️ CUSTOMIZE THIS SECTION ⚠️
    geometry = DomainGeometry("custom_geometry")
    # User adds geometry.add_domain() calls here
    
    # SECTION 4: PROBLEM CREATION (Automated from geometry)
    # SECTION 5: CONSTRAINT SETUP - ⚠️ CUSTOMIZE THIS SECTION ⚠️  
    # SECTION 6-7: GLOBAL DISCRETIZATION AND VALIDATION
\end{lstlisting}

\subsection{OrganOnChip Grid Problem}
\label{subsec:ooc_grid_problem}

The \texttt{ooc\_problem.py} implements a complete OrganOnChip problem with grid geometry.

\subsubsection{Main Framework Method}

\begin{lstlisting}[language=Python, caption=OoC Problem Framework]
def create_global_framework(geometry: Optional[DomainGeometry] = None,
                          config_file: Optional[str] = None):
    """
    OrganOnChip problem with grid geometry and configuration system.
    
    Args:
        geometry: Optional pre-defined geometry. Uses build_default_geometry() if None
        config_file: Optional TOML configuration file
    
    Returns:
        Tuple: (problems, global_discretization, constraint_manager, problem_name)
    """
    
    CONFIG_VALIDATION_TYPE = "ooc"      # For config file validation
    PROBLEM_TYPE = "organ_on_chip"      # For StaticCondensationFactory
    
    # Configuration loading with validation
    config_manager = OoCConfigManager()
    config = config_manager.load_config(config_file)
    
    # Extract configuration sections
    problem_config = config['problem']
    phys_params = config['physical_parameters']
    initial_conditions = config['initial_conditions']  # Already callables
    domain_initial_conditions = config.get('domain_initial_conditions', {})
    domain_force_functions = config.get('domain_force_functions', {})
    
    # Create problems from geometry
    problems = []
    for domain_id in range(geometry.num_domains()):
        domain_info = geometry.get_domain(domain_id)
        problem = Problem(neq, domain_info.domain_start, domain_info.domain_length,
                         parameters, "my_problem_type", name)
        # Set functions...
        problems.append(problem)
    
    # Setup constraints from geometry
    constraint_manager = setup_constraints_from_geometry(geometry, problems, neq)
    
    # Return framework components
    return problems, global_discretization, constraint_manager, problem_name
\end{lstlisting}

\subsubsection{Geometry Integration}

\begin{lstlisting}[language=Python, caption=Geometry Integration]
def build_default_geometry():
    """Build the default OoC grid geometry."""
    return build_grid_geometry()  # From geometry module

def setup_constraints_from_geometry(geometry: DomainGeometry, problems, neq: int):
    """Generate constraint manager from geometry connections."""
    constraint_manager = ConstraintManager()
    
    # Process boundary connections
    boundary_connections = geometry.get_boundary_connections()
    for conn in boundary_connections:
        if conn.is_exterior_boundary():
            for eq_idx in range(neq):
                constraint_manager.add_neumann(
                    equation_index=eq_idx,
                    domain_index=conn.domain1_id,
                    position=conn.parameter1,
                    data_function=lambda t: 0.0  # Homogeneous Neumann
                )
    
    # Process interior connections
    interior_connections = geometry.get_interior_connections()
    for conn in interior_connections:
        for eq_idx in range(neq):
            constraint_manager.add_trace_continuity(
                equation_index=eq_idx,
                domain1_index=conn.domain1_id,
                domain2_index=conn.domain2_id,
                position1=conn.parameter1,
                position2=conn.parameter2
            )
    
    return constraint_manager
\end{lstlisting}

\subsubsection{Domain-Specific Overrides}

\begin{lstlisting}[language=Python, caption=Domain-Specific Configuration]
# Apply domain-specific initial condition overrides
equation_names = ['u', 'omega', 'v', 'phi']

for domain_eq_key, func_name in domain_initial_conditions.items():
    # Parse "domain_<id>_<equation_name>"
    if domain_eq_key.startswith('domain_') and '_' in domain_eq_key[7:]:
        parts = domain_eq_key[7:].split('_', 1)
        domain_id = int(parts[0])
        equation_name = parts[1]
        
        if domain_id < len(problems) and equation_name in equation_names:
            equation_id = equation_names.index(equation_name)
            resolved_func = config_manager.function_resolver.resolve_function(func_name)
            problems[domain_id].set_initial_condition(equation_id, resolved_func)
            print(f"✓ Set initial condition: domain {domain_id}, equation '{equation_name}' -> {func_name}")
\end{lstlisting}

\subsection{Parameter Management}
\label{subsec:parameter_management}

\subsubsection{OrganOnChip Parameter Structure}

\begin{longtable}{|p{3cm}|p{2cm}|p{2cm}|p{6cm}|}
\hline
\textbf{Parameter} & \textbf{Index} & \textbf{Default} & \textbf{Physical Meaning} \\
\hline
\endhead

\texttt{nu} & 0 & 1.0 & Cell diffusion coefficient \\
\hline

\texttt{mu} & 1 & 2.0 & Potential diffusion coefficient \\
\hline

\texttt{epsilon} & 2 & 1.0 & Chemical diffusion coefficient \\
\hline

\texttt{sigma} & 3 & 1.0 & Velocity diffusion coefficient \\
\hline

\texttt{a} & 4 & 0.0 & Growth/decay rate \\
\hline

\texttt{b} & 5 & 1.0 & Velocity-potential coupling \\
\hline

\texttt{c} & 6 & 0.0 & Chemical reaction rate \\
\hline

\texttt{d} & 7 & 1.0 & Cell-chemical coupling \\
\hline

\texttt{chi} & 8 & 1.0 & Chemotactic sensitivity \\
\hline

\end{longtable}

\subsubsection{Keller-Segel Parameter Structure}

\begin{longtable}{|p{3cm}|p{2cm}|p{2cm}|p{6cm}|}
\hline
\textbf{Parameter} & \textbf{Config Path} & \textbf{Default} & \textbf{Physical Meaning} \\
\hline
\endhead

\texttt{mu} & diffusion.mu & 2.0 & Cell diffusion coefficient \\
\hline

\texttt{nu} & diffusion.nu & 1.0 & Chemical diffusion coefficient \\
\hline

\texttt{a} & reaction.a & 0.0 & Reaction parameter a \\
\hline

\texttt{b} & reaction.b & 1.0 & Reaction parameter b \\
\hline

\texttt{chi} & chemotaxis.chi & 'constant' & Chemotactic sensitivity function \\
\hline

\texttt{dchi} & chemotaxis.dchi & 'zeros' & Derivative of chemotaxis function \\
\hline

\end{longtable}

\subsection{Problem Module Usage Patterns}
\label{subsec:usage_patterns}

\subsubsection{Using OrganOnChip Problem}

\begin{lstlisting}[language=Python, caption=OoC Problem Usage]
# Method 1: Using default configuration
from bionetflux.problems.ooc_problem import create_global_framework
problems, global_disc, constraints, name = create_global_framework()

# Method 2: Using custom geometry
from bionetflux.geometry.domain_geometry import DomainGeometry
geometry = DomainGeometry("custom")
geometry.add_domain(extrema_start=(0,0), extrema_end=(1,0), name="domain1")
problems, global_disc, constraints, name = create_global_framework(geometry=geometry)

# Method 3: Using TOML configuration file
problems, global_disc, constraints, name = create_global_framework(
    config_file="my_ooc_config.toml"
)
\end{lstlisting}

\subsubsection{TOML Configuration Example}

\begin{lstlisting}[language=toml, caption=OoC TOML Configuration]
[problem]
name = "My_OoC_Problem"
neq = 4
problem_type = "ooc"

[time_parameters]
T = 2.0
dt = 0.05

[physical_parameters.viscosity]
nu = 1.5
mu = 2.5
epsilon = 1.2
sigma = 0.8

[physical_parameters.coupling]
chi = 2.0

[initial_conditions]
u = "ooc_gaussian_blob"
v = "ooc_parabolic_flow"

[domain_initial_conditions]
domain_0_u = "constant"
domain_3_v = "ooc_inlet_profile"

[domain_force_functions]
domain_1_omega = "sin_t"
\end{lstlisting}

\subsection{Function Resolution System}
\label{subsec:function_resolution}

\subsubsection{Function Library Hierarchy}

\begin{lstlisting}[language=Python, caption=Function Resolution Process]
# 1. Common functions (in BaseConfigManager)
common_functions = {
    'zeros': lambda s, t=0: np.zeros_like(s),
    'ones': lambda s, t=0: np.ones_like(s),
    'constant': lambda s, t=0: np.ones_like(s),
    'linear': lambda s, t=0: s,
    'quadratic': lambda s, t=0: s**2,
    'sin': lambda s, t=0: np.sin(np.pi * s),
    'cos': lambda s, t=0: np.cos(np.pi * s),
}

# 2. Problem-specific functions (OoC or KS)
# Added by _setup_function_library() in each config manager

# 3. Resolution order:
#    - Check problem-specific library first
#    - Fall back to common library
#    - Raise error if function not found
\end{lstlisting}

\subsection{Validation System}
\label{subsec:validation_system}

\subsubsection{Configuration Validation Rules}

\begin{lstlisting}[language=Python, caption=Validation Rules Example]
# In OoCConfigManager
def _setup_validation_rules(self):
    # Problem type validation
    self.validator.add_rule('problem.problem_type', 
                          {'type': str, 'required': True})
    
    # Physical parameter validation
    self.validator.add_rule('physical_parameters.viscosity.nu', 
                          {'type': float, 'min': 0.0, 'positive': True})
    
    # Time parameter validation
    self.validator.add_rule('time_parameters.T', 
                          {'type': float, 'min': 0.0, 'positive': True})
    
    # Discretization validation
    self.validator.add_rule('discretization.n_elements', 
                          {'type': int, 'min': 1})
\end{lstlisting}

\subsubsection{Problem Type Validation}

\begin{lstlisting}[language=Python, caption=Problem Type Validation]
def load_config(self, config_file: Optional[str] = None):
    config = super().load_config(config_file)
    
    # Validate problem type matches
    if config_file:
        config_problem_type = config.get('problem', {}).get('problem_type', None)
        if config_problem_type and config_problem_type != self.problem_type:
            raise ValueError(
                f"Configuration problem type '{config_problem_type}' does not match "
                f"expected problem type '{self.problem_type}'"
            )
    
    return config
\end{lstlisting}

\subsection{Mathematical Formulation}
\label{subsec:mathematical_formulation}

\subsubsection{OrganOnChip System}

The 4-equation OrganOnChip system implemented in the problems:

\begin{align}
\frac{\partial u}{\partial t} &= \nu \nabla^2 u - \chi \nabla \cdot (u \nabla \phi) + a u + f_u(x,t) \\
\frac{\partial \omega}{\partial t} &= \epsilon \nabla^2 \omega + c \omega + d u + f_\omega(x,t) \\
\frac{\partial v}{\partial t} &= \sigma \nabla^2 v + \lambda(\bar{\omega}) v + f_v(x,t) \\
\frac{\partial \phi}{\partial t} &= \mu \nabla^2 \phi + a \phi + b v + f_\phi(x,t)
\end{align}

where $\bar{\omega} = \frac{1}{|K|} \int_K \omega \, dx$ is the cell-averaged value.

\subsubsection{Keller-Segel System}

The 2-equation Keller-Segel system:

\begin{align}
\frac{\partial u}{\partial t} &= \mu \nabla^2 u - \nabla \cdot (u \chi(\phi) \nabla \phi) + f_u(x,t) \\
\frac{\partial \phi}{\partial t} &= \nu \nabla^2 \phi + u + f_\phi(x,t)
\end{align}

\subsection{Extension Guidelines}
\label{subsec:extension_guidelines}

\subsubsection{Creating New Problem Modules}

\begin{lstlisting}[language=Python, caption=New Problem Module Template]
def create_global_framework(config_file: Optional[str] = None):
    """
    New problem module following BioNetFlux patterns.
    """
    
    # 1. Create specialized config manager
    config_manager = MyNewConfigManager()
    config = config_manager.load_config(config_file)
    
    # 2. Extract parameters and functions
    parameters = extract_parameters_from_config(config)
    initial_conditions = config['initial_conditions']  # Already resolved
    
    # 3. Create geometry (default or provided)
    geometry = build_my_geometry()
    
    # 4. Create problems from geometry
    problems = []
    for domain_id in range(geometry.num_domains()):
        domain_info = geometry.get_domain(domain_id)
        problem = Problem(neq, domain_info.domain_start, domain_info.domain_length,
                         parameters, "my_problem_type", name)
        # Set functions...
        problems.append(problem)
    
    # 5. Setup constraints from geometry
    constraint_manager = setup_constraints_from_geometry(geometry, problems, neq)
    
    # 6. Return framework components
    return problems, global_discretization, constraint_manager, problem_name
\end{lstlisting}

\subsubsection{Configuration Manager Guidelines}

\begin{lstlisting}[language=Python, caption=New Config Manager Template]
class MyNewConfigManager(BaseConfigManager):
    def __init__(self):
        super().__init__("my_type")
    
    def _setup_defaults(self):
        self.defaults = {
            'problem': {'name': 'My_Problem', 'neq': N, 'problem_type': 'my_type'},
            'physical_parameters': { /* your parameters */ },
            'initial_conditions': { /* your variables */ },
            # ...existing code...
        }
    
    def _setup_function_library(self):
        my_functions = {
            'my_special_function': lambda s, t=0: /* implementation */,
            # ...existing code...
        }
        self.function_resolver.register_function_library(my_functions)
    
    def _setup_validation_rules(self):
        # Add your validation rules
        # ...existing code...
\end{lstlisting}

This documentation provides comprehensive coverage of the actual problem modules implemented in the BioNetFlux problems folder, reflecting the current configuration-driven architecture and providing clear guidelines for extension.

% End of problems folder detailed API documentation
